\chapter{Introduction}
According to wave-particle duality, light has both a particle and a wave character \cite{meschede2006gerthsen}.
In the wave picture, the wavelength, and thus also the frequency, determines whether the light is visible and what color it has.
For example, light with a wavelength between \SIrange{640}{430}{\nano\meter} is in the visible range.
Due to the electromagnetic properties of light, the entire spectrum is also referred to as the electromagnetic spectrum,
or EM spectrum for short.
This is shown in \autoref{fig:Spek}. \\
\begin{figure}[H]
    \centering
    \includegraphics[width=0.8\textwidth]{content/bilder/Em-Spek.png}
    \caption{Electromagnetic spectrum with frequency scale and typical fields of application.
    Marked in red is the THz range, which lies between the microwave and visible spectrum.
    Characteristic parameters of THz radiation are given below \cite{Williams_2006}.}
    \label{fig:Spek}
\end{figure}
While a large part of this spectrum can be generated without much effort, there are considerable challenges at certain frequency ranges.
This challenging range is the so-called terahertz gap; it describes
a frequency range in which it is technically particularly difficult to generate and detect electromagnetic waves.
This can be seen in \autoref{fig:THz-lucke}.
There, it can be seen that on the side of frequencies lower than THz, i.e. on the left side, the power of electronic generators
is plotted at different frequencies. On the right side, the powers of quantum cascade lasers at different frequencies
are plotted. The figure thus shows that just above \SI{1000}{GHz}, i.e. \SI{1}{THz}, the power on both sides
approaches 0. \\
This THz gap lies in the frequency range from \SIrange{0.1}{10}{\tera\hertz}, or, in the wavelength spectrum, from \SIrange{3000}{30}{\micro\meter} \cite{Williams_2006}.
Frequencies below the THz gap can be generated well using electronic oscillators.
However, as the frequency increases, more and more challenges arise, for example due to
parasitic capacitances, which result in less and less power being emitted as radiation.
In addition, simple black-body sources or diode lasers can generate high power
in the range above the THz gap.
However, this drops off sharply as the frequency decreases.
This results in a gap in the spectrum in which it is very difficult to generate an electromagnetic wave with THz frequencies
and high power.
\begin{figure}[H]
    \centering
    \includegraphics[width=0.8\textwidth]{content/bilder/THz-gap.png}
    \caption{Graph of the THz gap, where on the left of the gap the power of electrical radiation generators at lower frequencies is shown.
    On the right side, the power of quantum cascade lasers is plotted.
    Both powers decrease toward the marked THz gap in the middle \cite{thz_gap}.}
    \label{fig:THz-lucke}
\end{figure}
THz radiation is interesting for several reasons. Like X-rays, THz radiation can
be used to find and detect tumors and other tissue abnormalities, without causing the genetic damage
associated with X-rays \cite{pickwell2006biomedical}.
In the natural sciences, it is used for non-destructive material characterization and investigation,
as well as for low-energy excitation. \cite{PAWAR2013157}.
Nowadays, the generation of radiation within the THz gap is possible at higher intensities through the use of various THz generation techniques, such as photoconductive antennas
\cite{10.1117/1.OE.56.1.010901}, optical rectification, and THz quantum cascade lasers \cite{Dean_2014}.
This bachelor's thesis focuses in particular on optical rectification, since
this is the main mechanism by which the organic emitters generate THz radiation. \\ \\
They are compared with inorganic crystals. ZnTe is an inorganic crystal that can likewise generate THz radiation via optical rectification.
The crystal can be destroyed by excessively high optical power;
due to the small number of producers and the economic developments of recent years, the price has doubled over the last
five years. \\ \\
In this bachelor's thesis, two organic crystal systems are investigated with regard to their properties concerning THz generation and compared with
generation in ZnTe. \\
The organic crystals are grown/produced in the research group $look up which group this was$at TU Dortmund,
which enables rapid iteration regarding the design of the crystals. \\
This allows application-specific crystals to be grown with regard to size and thickness,
in contrast to commercially available inorganic crystals. \\
It is then discussed whether these crystals can replace the ZnTe crystals.
This thesis begins in chapter two with a theoretical description of the generation and detection of
THz radiation, building on which chapter three then explains terahertz time-domain spectroscopy.
This is abbreviated in what follows as THz-TDS, from the English term terahertz time-domain spectroscopy.
The results for the various crystals measured using this measurement method are presented and discussed in chapter four.
Chapter five concludes with an overview of the results and an outlook on further measurements.
